% from ega0-16-fr.tex
\oldpage[0\textsubscript{IV-1}]{38}
pour $M$, il résulte de (16.3.7) et (16.3.4) que $y_1$ n'appartient à aucun
des $\mathfrak{p}_j$, donc est $M$-régulier. Si on pose $P=M/y_1M$,
$(y_i)_{2\leq i\leq n}$ est un système de paramètres pour $P$, et il est
immédiat que $P$ vérifie l'hypothèse de l'énoncé~; appliquant l'hypothèse de
récurrence, on voit donc que $(y_i)_{2\leq i\leq n}$ est une suite
$P$-régulière, ce qui démontre la proposition.
\end{proof}

\begin{proposition}[16.5.9]
\label{0.16.5.9-fr}
Soient $A$ un anneau local noethérien, $M$ un $A$-module de type fini~; on
suppose que $M$ soit un module de Cohen-Macaulay. Soit
$\mathfrak{p}\in\operatorname{Supp}(M)$, et posons
$r=\dim(M)-\dim(M/\mathfrak{p}M)$. Alors il existe une suite $M$-régulière
$(x_i)_{1\leq i\leq r}$ formée d'éléments de $\mathfrak{p}$~; pour toute
suite ayant ces propriétés, on a
\[
  \dim(M/(x_1M+\cdots+x_rM))
  =\dim(M/\mathfrak{p}M)=\dim(A/\mathfrak{p}),
\]
et $\mathfrak{p}$ est un élément minimal de
$\operatorname{Ass}(M/(x_1M+\cdots+x_rM))$.
\end{proposition}

\begin{proof}
Prouvons d'abord l'existence des $x_i$~; il n'y a rien à démontrer pour
$r=0$. Pour $r>0$, procédons par récurrence sur $r$. Comme pour tout idéal
premier $\mathfrak{q}_i$ associé à $M$ on a
$\dim(A/\mathfrak{q}_i)=\dim(M)$ (16.5.4), l'hypothèse
$\dim(M/\mathfrak{p}M)<\dim(M)$ entraîne que $\mathfrak{p}$ ne peut être
contenu dans un des $\mathfrak{q}_i$, ces derniers étant les idéaux premiers
minimaux parmi ceux qui contiennent $\operatorname{Ann}(M)$ (16.1.2.2 et
16.1.7)~; on en conclut que $\mathfrak{p}$ n'est pas non plus contenu dans la
réunion des $\mathfrak{q}_i$ (Bourbaki, \emph{Alg.~comm.}, chap.~II, \S~1,
no~1, prop.~2), et par suite il existe un $x_1\in\mathfrak{p}$ qui est
$M$-régulier (Bourbaki, \emph{Alg.~comm.}, chap.~IV, \S~1, no~1, cor.~2 de
la prop.~2). On déduit donc de (16.5.5) que $N=M/x_1M$ est un module de
Cohen-Macaulay, de dimension $\dim(M)-1$~; comme $x_1\in\mathfrak{p}$, on a
d'ailleurs $N/\mathfrak{p}N=M/\mathfrak{p}M$, d'où
$\dim(N)-\dim(N/\mathfrak{p}N)=r-1$~; on peut par suite appliquer à $N$
l'hypothèse de récurrence, et si $(x_i)_{2\leq i\leq r}$ est une suite
$N$-régulière formée d'éléments de $\mathfrak{p}$, il est clair que
$(x_i)_{1\leq i\leq r}$ est une suite $M$-régulière formée d'éléments de
$\mathfrak{p}$.

Si on pose $P=M/(x_1M+\cdots+x_rM)$, on a
$P/\mathfrak{p}P=M/\mathfrak{p}M$ puisque les $x_i$ sont dans
$\mathfrak{p}$, et $\dim(P)=\dim(P/\mathfrak{p}P)$ par (16.5.6)~; mais comme
$P$ est un module de Cohen-Macaulay, on a aussi
$\dim(P)=\dim(A/\mathfrak{p}')$ pour tout
$\mathfrak{p}'\in\operatorname{Ass}(P)$ (16.5.4). Comme
$\mathfrak{p}\in\operatorname{Supp}(P)$ ($0_{\mathrm I}$, 1.7.5),
$\mathfrak{p}$ contient un des idéaux
$\mathfrak{p}'\in\operatorname{Ass}(P)$ (Bourbaki, \emph{Alg.~comm.},
chap.~IV, \S~1, no~4, th.~2), et comme
\[
  \dim(P)=\dim(P/\mathfrak{p}P)
  \leq\dim(A/\mathfrak{p})
  \leq\dim(A/\mathfrak{p}')=\dim(P),
\]
on a nécessairement $\mathfrak{p}=\mathfrak{p}'$ (16.1.2.2), ce qui achève
la démonstration.
\end{proof}

\begin{corollary}[16.5.10]
\label{0.16.5.10-fr}
Sous les mêmes hypothèses que dans (16.5.9)~:
\begin{enumerate}
  \item[(i)] $M_{\mathfrak{p}}$ est un $A_{\mathfrak{p}}$-module de
  Cohen-Macaulay.
  \item[(ii)] On a
  \begin{equation}
  \label{0.16.5.10.1-fr}
  \tag{16.5.10.1}
    \dim(M)=\dim_A(M/\mathfrak{p}M)
    +\dim_{A_{\mathfrak{p}}}(M_{\mathfrak{p}}).
  \end{equation}
\end{enumerate}
\end{corollary}

\begin{proof}
Avec les notations de (16.5.9), soient $x_i'$ les images canoniques des $x_i$
dans $A_{\mathfrak{p}}$ $(1\leq i\leq r)$~; comme les $x_i'$ appartiennent à
l'idéal maximal de $A_{\mathfrak{p}}$ et forment une suite
$M_{\mathfrak{p}}$-régulière par platitude (15.1.14), on a
\[
  \operatorname{prof}_{A_{\mathfrak{p}}}(M_{\mathfrak{p}})
  \geq r=\dim(M)-\dim(M/\mathfrak{p}M)
  \geq\dim_{A_{\mathfrak{p}}}(M_{\mathfrak{p}})
\]
(en vertu de (16.1.8.1)). Mais compte tenu de (16.4.5.1), les trois termes de
cette inégalité sont nécessairement égaux, d'où le corollaire.
\end{proof}

